\subsection{Personalized Fast FullSubNet}

\subsubsection{Motivation}

The Fast FullSubNet (FFSN) model \cite{hao2023fast} offers a lightweight and efficient subband-based speech enhancement model that maintains high performance in noise reduction tasks.
However, the model is trained on a large dataset with a wide range of noise types and SNR levels, which may not be optimal for personalized speech enhancement.
To introduce personalization, we propose a personalized Fast FullSubNet (PFFSN) model leveraging the model ability to process subbands. PFFSN consists of two major modules: a speaker embedding module and a denoising module.
The speaker embedding module is responsible for extracting speaker-specific features from the input speech signal, while the denoising module is responsible for enhancing the speech signal based on the extracted speaker features.
The architecture of the PFFSN model is shown in Figure \ref{fig:pffsn-architecture}.

\subsubsection{Speaker Embedding Module}
The speaker embedding module consists of the encoder and bottleneck layers of FFSN followed by a simple feedforward network.
The encoder and bottleneck layers are responsible for extracting subband features of the speaker given a clean speech signal.
The feedforward network aims to classify whether the targeted speaker voice resides in a subband or not.
The feedforward network consists of two fully connected layers with LayerNorms and ReLU activation functions and a sigmoid activation function in the output layer. 
The output of the sigmoid activation is the probability of the speaker voice residing in the subband.

\subsubsection{Denoising Module}
The denoising module is a regular FFSN model in which its decoder inputs subbands are scaled based on the speaker embedding module's output.
Subbands in which the speaker voice is detected are amplified, while subbands in which the speaker voice is not detected are suppressed.
This process aims to help the decoder focus on enhancing the speaker voice while suppressing the noise.


\subsubsection{Experimental Setup}
We trained the PFFSN model on the training set of our dataset for 43 epochs with a batch size of 16 per GPU. We trained the model using Kaggle's 2 NVIDIA Tesla T4 GPUs.
The Adam optimizer was used with a learning rate of $5 \times 10^{-5}$ with Cosine Decay and beta values of 0.9 and 0.999. Model exponential moving average (EMA) was used with a decay of 0.999.
We employed a mixture of loss functions, including L1 loss, Multi Resolution STFT loss, complex Ration Mask (cRM) loss and TAP loss with weights 1, 0.5, 1 and 0.3 respectively.